% Additive bridge between the two received calculations. Dependencies are
% checked in the adjacent transport/equilibrium/asymptotics/tail audits.
\begingroup
\section{The exact profile--covariance dictionary and quantitative joint return}
\label{sec:PCB}

All objects below retain the original simple quartet, its period, its
conductor coefficients, and the two source orders $s=1,a$. Put
$a\equiv1\pmod4$, $a\ge257$, $q=(a+1)^2$, $Q=(a-7)^2$,
$\Delta=q-Q=16a-48$, $m=\Delta-v$, and $\ell_s=(s-1)/4$.
The original full polynomial and the attained old-relation minimum are
those of NTG; none is replaced by a diagonal norm. This section connects
the two web derivations received on 18 September 2026. Their original
statements and their reviewed corrections are retained in the accompanying
source records. The equilibrium identities used below were proved for
these exact conventions in EIQ, ECL6, ECL14--17, and GRE6--8.

Every finite comparison below is on the intersection of the two original
root-comparison domains. Explicitly, retain $0<\delta<1/2$, $\gamma>2$,
$v=\operatorname{ord}_0E_A\in\{0,\ldots,80\}$, $\mu_v\ne0$, and impose
\[
 a\sqrt{\delta^2+\gamma^2}\le2^{-33}q,\qquad
 n_b=(b+1)^2+\ell_s\ge100,\quad
 \max\{b\sqrt{\delta^2+\gamma^2},2\ell_s+1/2\}
       \le2^{-17}n_b\quad(b=a,a-8).
 \tag{PCB0}
\]
These are the stated attained-envelope and retained-Gamma covariance
guards. They hold eventually for every fixed actual quartet and both
orders $s=1,a$. In this domain $Q>0$, $m=\Delta-v>0$ and $0<R<q$;
the lower row cutoffs obey the original whole-polynomial degree bounds.
The scalar identities PCB1--6 hold for every $t>0$ independently of
these finite-family guards.

\subsection{An exact dictionary of the two scalar functions}

For $t>0$, let $u(t),v_*(t)$ be the endpoints for
$\alpha=2/t$, $\beta=\pi/t$, and set
\[
 w=(u+v_*)/2,\qquad c=(v_*^2-u^2)/4,\qquad
 \rho=u/v_*,\qquad k=(1-\rho^2)^{1/2}.
 \tag{PCB1}
\]
The actual EIQ equations give $uK(k)=2$ and
$v_*E(k)=2(1+t)$. In particular
$\rho K(k)/E(k)=1/(1+t)$; this is precisely the implicit parameter
in the attained-tail response. Retain the actual Robin constant
\[
 \lambda=2(\alpha+2)-2\alpha\log w-2\log c,
 \qquad \lambda_\alpha=-2\log w,
 \qquad \lambda_\beta=2(\alpha+2)/\beta.
 \tag{PCB2}
\]
These are GRE6--8, including the endpoint motion and the original
Robin sign. Thus $\lambda-\alpha\lambda_\alpha-
\beta\lambda_\beta=-2\log c$. Differentiating the explicit function
\[
 \psi(t)=(1+t)(1-\log(1+t))
       -\tfrac t4\lambda(2/t,\pi/t)
 \tag{PCB3}
\]
by the chain rule proves
\[
 \psi'(t)=-\log(1+t)+\tfrac12\log c
         =\log(k/E(k)),
 \qquad
 2\psi(t)-2(1+t)\psi'(t)
    =\log\frac{w^2}{c}
    =\log\frac{1+\rho}{1-\rho}=:\mathfrak f(t).
 \tag{PCB4}
\]
For the middle equality substitute PCB2: the terms $2(1+t)$ cancel
$t(\alpha+2)$ exactly because $t\alpha=2$. The remaining expression
is $2\log w-\log c$. This proves the exact map between the covariance
function and the attained-row profile, rather than identifying them by
their integrated constants.

Write $a_0=\psi(0)=\log(4/\pi)$, $a_1=\psi(1)$, and
$C_B=4\int_0^1\psi(t)\,dt$. The convergent elliptic-integral expansions
at $k=0$, or the explicit expansion in the covariance proof, give
\[
 \psi(t)=a_0+\tfrac14t\log t+(a_0-\tfrac14)t+O(t^{3/2}),
 \quad \mathfrak f(t)=-\tfrac12\log t+O(\sqrt t).
 \tag{PCB5}
\]
In particular the improper integral of $\mathfrak f$ exists. Integrating
PCB4 on $[\epsilon,1]$ and then sending $\epsilon$ to zero gives
\[
 2\int_0^1\mathfrak f(t)\,dt
 =8\int_0^1\psi(t)\,dt-8a_1+4a_0
 =2C_B-8a_1+4a_0=C_\partial.
 \tag{PCB6}
\]
The final equality is the original ECL16 coefficient expressed by its
Robin/energy primitive. All constants, including the factor two for the
two-endpoint return, are therefore the same in the two derivations.

\subsection{A quantitative improvement obtained by using both results}

The native row is
$\gamma_{r,s}=\log(\eta_{r,s}/\nu^-_{m+r,s})\ge0$, where $\eta$
is the complete attained preceding-relation minimum. Its original weights
are $\vartheta_0=\vartheta_q=1$ and $\vartheta_r=2$ for $0<r<q$.
Set $R=\lfloor a^{3/2}/\log a\rfloor$ and $B=\{R+1,\ldots,q\}$.
Every sum in this subsection has this finite range, not the larger range
available to the envelope for the proper-source displacement.

The audited covariance row estimate and the original finite endpoint
factorization give, uniformly in this range,
\[
 \left|\sum_{r\in B}\vartheta_r\gamma_{r,s}
  -\sum_{r\in B}\vartheta_r
    \{G(q+\ell_s,r)-G(Q+\ell_s,m+r)\}\right|
 \le E^{\rm nat}_{a,s}(R)=O_{\rm actual}(q),
 \quad G(n,r)=2n\psi(r/n).
 \tag{PCB7}
\]
Here $E^{\rm nat}$ is a finite certified allowance: sum the absolute
row errors from the covariance norm intervals over the two grids, and add
the absolute endpoint allowances for the conductor and analytic quotient
at $q+R,2q-1,2q$, together with the low-factor tail bound. Those row errors
are bounded by
$C_{\rm actual}\{1+\log^+(n/(r+1))+
(\ell_s+1)^2/(n+r)\}$; the lower row uses $m+r$.
Their sums are $O(q)$. The transport endpoint allowances are $O(q)$
at every rank, and the low-factor tail is between zero and its proved
$O(q)$ full return. This proves PCB7 by the exact TAC factorization,
with no replacement of the actual tail by independent endpoint choices.

There is a direct finite-difference identity behind the leading profile:
\[
 (\partial_n-\partial_r)G(n,r)=\mathfrak f(r/n),\qquad
 G(q+\ell_s,r)-G(Q+\ell_s,r+\Delta)
 =\int_{Q+\ell_s}^{q+\ell_s}
  \mathfrak f\!\left(\frac{q+\ell_s+r-x}{x}\right)\,dx.
 \tag{PCB8}
\]
It follows by differentiating $G(x,q+\ell_s+r-x)$ and using PCB4.
The function $t\mathfrak f'(t)$ is bounded on $0<t\le4$:
away from zero this follows from the nonzero derivative in the EIQ
implicit endpoint equation; at zero the convergent expansion in powers
of $\sqrt t$ in PCB5 gives $t\mathfrak f'(t)=-1/2+O(\sqrt t)$.
Also $|\psi'(t)|\le C(1+|\log t|)$ on that interval by PCB4--5.
For $x\in[Q+\ell_s,q+\ell_s]$, the logarithmic difference of the
argument in PCB8 from $r/q$ is bounded by
$C\{\Delta/r+(\Delta+\ell_s)/q\}$. The mean-value integral for
$\mathfrak f$ with respect to $\log t$ now proves
\[
 \left|G(q+\ell_s,r)-G(Q+\ell_s,m+r)
                 -\Delta\mathfrak f(r/q)\right|
 \le C_{\rm actual}(q/r+\log q),\qquad r\in B.
 \tag{PCB9}
\]
The difference between $m+r$ and $r+\Delta$ is exactly the fixed integer
$v$ and costs at most $2v\sup|\psi'|=O_{\rm actual}(\log q)$.
The other terms use $\Delta^2=O(q)$ and $\Delta\ell_s=O(q)$.
Summing PCB9 gives a finite allowance
$E^{\rm cmp}_{a,s}(R)=O_{\rm actual}(q\log q)$.

Let $U_{a,s,r}$ be the original un-intersected conductor/optimized-trial
upper endpoint in the attained-tail proof, with its full conductor factor,
original mass cancellation and lower-ideal minimum. A smaller endpoint
obtained by intersecting with NIB may bound the attained row but need not
bound this particular trial; it is not used for the trial-slack statements.
Put
\[
 E^U_{a,s}(R)=\sum_{r\in B}\vartheta_r
       |U_{a,s,r}-\Delta\mathfrak f(r/q)|,
 \qquad E^T_{a,s}(R)=E^{\rm nat}_{a,s}(R)+E^{\rm cmp}_{a,s}(R).
 \tag{PCB10}
\]
The finite envelope gives $0\le\gamma\le U$ and
$E^U=O_{\rm actual}(q\log q)$. Combining PCB7--10 gives
\[
 \begin{split}
 0&\le\sum_{r\in B}\vartheta_r(U_{a,s,r}-\gamma_{r,s})
       \le E^U_{a,s}(R)+E^T_{a,s}(R),\\
 \sum_{r\in B}\vartheta_r
       |\gamma_{r,s}-\Delta\mathfrak f(r/q)|
       &\le2E^U_{a,s}(R)+E^T_{a,s}(R)
        =O_{\rm actual}(q\log q).
 \end{split}\tag{PCB11}
\]
Indeed subtract the actual tail sum from the upper-endpoint sum for the
first inequality. For the second use pointwise
$|\gamma-h|\le(U-\gamma)+|U-h|$, valid whenever $\gamma\le U$.
This uses the actual one-sided envelope, not a pointwise lower estimate.
If $\tau_{r,s}$ denotes the optimized upper-fibre trial norm, the exact
projection identity from the attained-tail proof further gives
\[
 0\le\sum_{r\in B}\vartheta_r\log(\tau_{r,s}/\eta_{r,s})
       \le E^U_{a,s}(R)+E^T_{a,s}(R)
       =O_{\rm actual}(q\log q).
 \tag{PCB12}
\]
Thus the two independently derived results together replace the earlier
unquantified $o(aq)$ tail-profile and trial-slack conclusion by a proved
$O(q\log q)$ allowance on the unchanged retained tail. Relative to $aq$
this is $O_{\rm actual}(\log a/a)$. It remains on the scale relevant to
the final exterior benchmark; no sign at that scale has been inferred.

The estimate receives every original subinterval of $B$, since restricting
an absolute-error sum can only decrease it. It gives no individual
allocation coefficient: the exact identities still read
$S_K+S_R=\mathcal T$ and
$4S_K=\mathcal L_e+V_b+V_{\rm rem}+V_s+V_A$.
The independent degree-$k$ proper source, original arithmetic correction
and projected pairing remain in their own proved receiving maps.

\subsection{The full row range and a uniform rate for actual mixed determinants}

The attained-envelope proof also applies for every integer
$2\le r\le q+32a$. Its finite norm estimate uses
$h=\lfloor r/2\rfloor\ge1$, and the error $O(q/r+\log q)$ explicitly
permits $r$ fixed. No expansion in a small $\Delta/r$ is used: in PCB8,
$\log(1+x)\le x$ bounds the logarithmic argument change for all $x\ge0$.
The parity shifts in the exact Robin parameters cost $O(q/r)$, and the
original root comparison holds before minimization over the full stated
polynomial range. These details are proved in the attained-envelope
section; they extend its initially stated tail range without changing its
objects.

Retain row zero as its actual finite atom $\gamma_{0,s}$, rather than
assigning a value to $\mathfrak f(0)$. The existing NIB estimates give
\[
 0\le\gamma_{0,s},\gamma_{1,s}
       \le C_{\rm actual}a\log q,
 \qquad h_r:=\Delta\mathfrak f(r/q)\quad(1\le r\le q).
 \tag{PCB13}
\]
Because $\mathfrak f$ is positive and decreasing and
$\mathfrak f(t)\le C(1+\log(1/t))$ on $(0,1]$, comparison on each mesh
interval, including the first improper integral, gives
\[
 \left|\sum_{r=1}^{q}\vartheta_rh_r-\Delta qC_\partial\right|
       \le C\Delta\log q.
 \tag{PCB14}
\]
For clarity, the lower rectangle sum $q^{-1}\sum_{r=1}^q\mathfrak f(r/q)$
is below the integral. All discrepancies after the first cell telescope
to at most $\mathfrak f(1/q)/q$, and the first cell costs
$\int_0^{1/q}\mathfrak f(t)dt=O(\log q/q)$.
Multiplication by $2\Delta q$, followed by subtraction of the single
extra endpoint value $\Delta\mathfrak f(1)$, proves PCB14.

The now evaluated native total, with its original sign, is
\[
 \mathcal T_{a,s}=16C_\partial aq-128q\log a+O_{\rm actual}(q)
       =\Delta qC_\partial+O_{\rm actual}(q\log q).
 \tag{PCB15}
\]
The second equality uses $\Delta=16a-48$; the omitted term is exactly
$48C_\partial q$ together with the displayed $-128q\log a$, and no
favorable sign is selected. PCB13--15 imply
$|\sum_{r=2}^q\vartheta_r(\gamma_{r,s}-h_r)|=O(q\log q)$.
The all-row envelope has
$\sum_{r=2}^q\vartheta_r|U_{a,s,r}-h_r|=O(q\log q)$, since
$\sum_{r=2}^q q/r=O(q\log q)$.
Apply the two elementary inequalities used in PCB11 to this whole finite
range, and then restore row one using PCB13. The result is
\[
 \boxed{\sum_{r=1}^q\vartheta_r
        |\gamma_{r,s}-\Delta\mathfrak f(r/q)|
          =O_{\rm actual}(q\log q),\qquad
 0\le\gamma_{0,s}=O_{\rm actual}(a\log q).}
 \tag{PCB16}
\]
The optimized-trial slack has the same $O(q\log q)$ bound over
$2\le r\le q$. No trial-slack assertion for rows zero or one is inferred
from NIB, which bounds the attained row rather than that particular trial.

Let $0\le p\le r\le q$ be any original indices and
$Y=[x_{p+1},\ldots,x_r]$. The complete NTG determinant identity gives
\[
 \log\det(I+Y^*\Omega_p^{-1}Y)
      =\sum_{j=p+1}^{r}\gamma_{j,s}.
\]
Restricting the absolute sum in PCB16 bounds the difference from
$\Delta\sum_{j=p+1}^r\mathfrak f(j/q)$ by $O(q\log q)$,
uniformly in both indices. The same monotone integral comparison as PCB14
costs at most $O(\Delta\log q)$ uniformly, even if $p=0$.
Since $\Delta/a=16-48/a$, we obtain
\[
 \boxed{\sup_{0\le p\le r\le q}
 \left|\frac{\log\det(I+Y^*\Omega_p^{-1}Y)}{aq}
       -16\int_{p/q}^{r/q}\mathfrak f(t)dt\right|
       =O_{\rm actual}(\log a/a).}
 \tag{PCB17}
\]
Every old inverse and every later cross pairing in this determinant is
retained. The two-endpoint tail has the corresponding coefficient $32$;
the last single row costs $O(q\log q)$ by PCB16 and so obeys the same
uniform rate. These are quantitative versions of the actual mixed
determinant profile, not pointwise lower estimates on every row.

Finally, for any set $I\subset\{0,\ldots,q\}$ of cardinality $d\ge1$,
decreasingness of $\mathfrak f$ places its largest $d$ positive-index
values first. The bound in PCB14's proof or
$\sum_{j=1}^d\log(q/j)\le d\{1+\log(q/d)\}$ therefore gives
\[
 0\le\mathcal T^I_{a,s}
    \le C_{\rm actual}\{ad(1+\log(q/d))+q\log q\}.
 \tag{PCB18}
\]
For $d=q+1$ use the same statement with $d$ replaced by $q$, plus the
row-zero allowance in PCB13. For the empty set the return is exactly zero.
All original two allocations and four CPQ complements receive this bound
with their established coefficients one and four, respectively, through
the existing exact flag maps. In particular every $d=o(q)$ set is
$o(aq)$. The signed allocation or projected arithmetic pairing still
requires its own calculation; PCB18 supplies no replacement sign.
\endgroup
